% ============================================================
% PAGE 8 — Experimental Setup, Datasets, and Training Details
% ============================================================

\section{Experimental Results and Discussion}
\label{sec:results}

This section details the experimental protocol used to evaluate
FusionDet, beginning with the benchmark datasets, proceeding
through the hardware and software configuration, and concluding
with the formal definition of every evaluation metric reported
in subsequent analyses.  Reproducibility is a first-order
concern: every hyperparameter, library version, and random seed
is recorded so that independent groups can replicate our
results without ambiguity.

\subsection{Datasets}

\subsubsection{MS-COCO 2017}

The Microsoft Common Objects in Context (MS-COCO)
\cite{lin2014microsoft} dataset is the primary benchmark for
all experiments.  COCO 2017 comprises 118\,287 training images
(\texttt{train2017}), 5\,000 validation images
(\texttt{val2017}), and 20\,288 test images
(\texttt{test-dev2017}) spanning 80 object categories.  Its
defining characteristic is the diversity of object scales: the
evaluation protocol partitions detections into \textit{small}
($\text{area} < 32^{2}$~px), \textit{medium}
($32^{2} \leq \text{area} < 96^{2}$~px), and \textit{large}
($\text{area} \geq 96^{2}$~px) subsets, providing granular
visibility into a detector's scale-dependent performance.
All primary results are reported on \texttt{val2017}; final
comparative numbers are additionally submitted to the
\texttt{test-dev} evaluation server for independent
verification.

\subsubsection{PASCAL VOC 2012}

As a secondary benchmark we employ PASCAL VOC
2012~\cite{everingham2010pascal}, which contains 11\,540
\texttt{trainval} images across 20 categories.  Although
smaller and less challenging than COCO, VOC remains widely used
in the literature and provides a useful cross-dataset
generalization check.  We train on the combined VOC 2007+2012
\texttt{trainval} split (16\,551 images) and evaluate on the
VOC 2007 \texttt{test} split (4\,952 images).

\subsubsection{Data Augmentation Pipeline}

Both datasets are processed through an identical augmentation
pipeline: mosaic augmentation~\cite{bochkovskiy2020yolov4}
(four images composited into one), random horizontal flip
($p = 0.5$), colour jitter (brightness $\pm 0.4$, contrast
$\pm 0.4$, saturation $\pm 0.7$, hue $\pm 0.015$), random
affine transformations (rotation $\pm 10^{\circ}$, scale
$0.1$--$2.0$, translation $\pm 0.1$), and MixUp
\cite{zhang2018mixup} with $\alpha = 1.5$ applied
stochastically during the first 280 epochs.  Mosaic and MixUp
are disabled for the final 20 epochs to allow batch-normalization
statistics to stabilize on unperturbed data~\cite{ge2021yolox}.

\subsection{Hardware and Software Configuration}

Table~\ref{tab:hw_sw} summarizes the complete training
environment.

\begin{table}[!t]
\centering
\caption{Hardware and software configuration.}
\label{tab:hw_sw}
\renewcommand{\arraystretch}{1.15}
\setlength{\tabcolsep}{4pt}
\begin{tabular}{l l}
\hline\hline
\textbf{Component}      & \textbf{Specification} \\
\hline
GPU                      & 4$\times$ NVIDIA RTX 4090 (24\,GB) \\
CPU                      & AMD EPYC 7543 (32 cores) \\
RAM                      & 256\,GB DDR4-3200 \\
Framework                & PyTorch 2.1.0, CUDA 12.1 \\
Mixed Precision          & FP16 via \texttt{torch.cuda.amp} \\
Distributed Strategy     & PyTorch DDP (1 node, 4 GPUs) \\
\hline
Optimizer                & AdamW~\cite{loshchilov2019decoupled} \\
Base Learning Rate       & $1 \times 10^{-3}$ \\
Weight Decay             & $5 \times 10^{-2}$ \\
LR Schedule              & Cosine annealing to $1 \times 10^{-5}$ \\
Warm-up                  & Linear, 3 epochs \\
Batch Size               & 64 (16 per GPU) \\
Input Resolution         & $640 \times 640$ \\
Total Epochs             & 300 \\
EMA Decay                & 0.9999 \\
Random Seed              & 42 \\
\hline\hline
\end{tabular}
\end{table}

The AdamW optimizer~\cite{loshchilov2019decoupled} is chosen
over vanilla SGD for its decoupled weight-decay formulation,
which improves generalization in transformer-containing
architectures by preventing interaction between adaptive
learning-rate scaling and the regularization term.  The cosine
annealing schedule smoothly reduces the learning rate from
$1 \times 10^{-3}$ to $1 \times 10^{-5}$, avoiding the abrupt
drops associated with step-decay schedules.  An exponential
moving average (EMA) of model weights with decay $0.9999$ is
maintained throughout training and used for all evaluation.

Inference latency is measured on a single RTX~4090 GPU at
batch size 1, FP16 precision, and \texttt{torch.no\_grad()},
averaged over the full \texttt{val2017} split after a 50-image
warm-up.

\subsection{Evaluation Metrics}

\subsubsection{Mean Average Precision (mAP)}

For each class $k$, a precision--recall curve is constructed
by sweeping a confidence threshold.  The Average Precision
(AP) is the area under the interpolated curve:
\begin{equation}
    \text{AP}_{k}^{\tau} =
    \int_{0}^{1} p_{k}^{\tau}(r)\; dr
    \label{eq:ap}
\end{equation}
The mean across all $K$ classes:
\begin{equation}
    \text{mAP}^{\tau} =
    \frac{1}{K}\sum_{k=1}^{K} \text{AP}_{k}^{\tau}
    \label{eq:map_tau}
\end{equation}
COCO further averages over ten IoU thresholds from $0.50$ to
$0.95$ in steps of $0.05$:
\begin{equation}
    \text{mAP}@[.50{:}.05{:}.95] =
    \frac{1}{10}\sum_{t=1}^{10}
    \text{mAP}^{0.45 + 0.05\,t}
    \label{eq:map_coco}
\end{equation}
We report this metric together with $\text{mAP}@0.50$,
$\text{AP}_{S}$, $\text{AP}_{M}$, and $\text{AP}_{L}$.

\subsubsection{Inference Speed (FPS)}

Frames Per Second is computed as:
\begin{equation}
    \text{FPS} = \frac{N_{\text{images}}}
    {\sum_{i=1}^{N_{\text{images}}} t_{i}}
    \label{eq:fps}
\end{equation}
where $t_{i}$ is the wall-clock time for image $i$, inclusive
of pre-processing (resize, normalize) and post-processing (NMS).

\subsubsection{Computational Complexity}

We additionally report the number of floating-point operations
(GFLOPs) and trainable parameters~(M) for every model variant,
measured at reference resolution $640 \times 640$ using the
\texttt{fvcore} \texttt{FlopCountAnalysis} utility to ensure
consistency with values reported in the literature.
